\section{Limits of functions}

A function limit describes what happens to the values \(f(x)\) when the input \(x\) approaches a certain point. The point may or may not belong to the domain of the function. What matters for the limit is the behavior near the point.

We write
\[
\lim_{x\to a} f(x)=L
\]
when the values \(f(x)\) approach \(L\) as \(x\) approaches \(a\).

\begin{definitionbox}
The notation
\[
\lim_{x\to a} f(x)=L
\]
means that the values of \(f(x)\) can be made arbitrarily close to \(L\) by taking \(x\) sufficiently close to \(a\), with \(x\ne a\).
\end{definitionbox}

The condition \(x\ne a\) is not a technicality. A limit studies nearby values. The value at the point itself is a separate question.

\subsection*{Direct substitution}

Sometimes the limit can be found by direct substitution. This often happens for polynomials and many simple functions.

\begin{examplebox}
Compute
\[
\lim_{x\to3}(x^2+2x-1).
\]
The expression is defined and behaves regularly near \(x=3\), so we substitute:
\[
3^2+2\cdot3-1=9+6-1=14.
\]
Therefore
\[
\lim_{x\to3}(x^2+2x-1)=14.
\]
\end{examplebox}

Direct substitution is not a rule to use blindly. It is a first test. If it gives a meaningful value and the function behaves regularly, it may solve the problem. If it produces a form such as \(0/0\), more work is needed.

\subsection*{A removable difficulty}

Consider
\[
\lim_{x\to3}\frac{x^2-4x+3}{x-3}.
\]
Direct substitution gives
\[
\frac{9-12+3}{0}=\frac0{0},
\]
which is not a number. This form tells us that direct substitution has failed. It does not tell us that the limit does not exist.

Factor the numerator:
\[
x^2-4x+3=(x-1)(x-3).
\]
For \(x\ne3\),
\[
\frac{x^2-4x+3}{x-3}=x-1.
\]
Therefore
\[
\lim_{x\to3}\frac{x^2-4x+3}{x-3}
=\lim_{x\to3}(x-1)=2.
\]

The original expression is not defined at \(x=3\), but near \(3\) it behaves like \(x-1\). The limit reads the nearby behavior.

\subsection*{The meaning of \(0/0\)}

The expression \(0/0\) is called an indeterminate form. It means that the information from direct substitution is insufficient. It is not an answer.

The limit may exist, may fail to exist, or may be infinite. We must transform the expression or analyze it in another way.

\begin{examplebox}
The limit
\[
\lim_{x\to1}\frac{x^2-1}{x-1}
\]
gives \(0/0\) by direct substitution. After factoring,
\[
\frac{x^2-1}{x-1}=x+1
\]
for \(x\ne1\), so the limit is \(2\).

But
\[
\lim_{x\to0}\frac{|x|}{x}
\]
also cannot be found by direct substitution, and the two-sided limit does not exist because the left and right behaviors are different.
\end{examplebox}

This is why we must read the situation before choosing a technique.

\subsection*{Limits and graphs}

On a graph, the limit \(\lim_{x\to a} f(x)\) asks what height the graph approaches near \(x=a\). A hole in the graph may not prevent a limit from existing. A jump may prevent a two-sided limit from existing.

\begin{remarkbox}
The value \(f(a)\) is the height of the graph at \(x=a\), if that point exists. The limit \(\lim_{x\to a}f(x)\) is the height approached by nearby points of the graph.
\end{remarkbox}

\begin{summarybox}
When computing a function limit, ask:
\begin{itemize}
\item What point is \(x\) approaching?
\item Is the function defined at nearby points?
\item Does direct substitution give a meaningful value?
\item If direct substitution gives \(0/0\), can the expression be simplified?
\item Is the left-hand behavior the same as the right-hand behavior?
\item Does the graph suggest a hole, a jump, or an asymptote?
\end{itemize}
\end{summarybox}

A function limit is a statement about local behavior. The point itself matters for continuity, but the limit first asks what happens nearby.